\appendix
\section{Mathematical Formulation of Sacroiliac Joint Micromotion Analysis}
\label{app:mathematical_formulation}

This appendix provides a self-contained, mathematically rigorous specification of the computational pipeline used to extract three-dimensional relative sacroiliac joint (SIJ) kinematics, boundary conditions, and morphometric dimensions from the finite-element (FE) cohort simulations.

\subsection{Continuum finite-element formulation and boundary value problem}
The mechanical equilibrium of the human pelvic ring is governed by the weak form of small-strain linear elasticity:
\begin{equation}
\int_{\Omega} \boldsymbol{\sigma}(\mathbf{u}) : \boldsymbol{\varepsilon}(\mathbf{v}) \, d\Omega + \int_{\Gamma_{S1}} \gamma \, (\mathbf{u} \cdot \mathbf{v}) \, ds = \int_{\Gamma_{\text{load}}} \mathbf{t}_{\text{ext}} \cdot \mathbf{v} \, ds, \quad \forall \mathbf{v} \in \mathcal{V},
\label{eq:app_weak_form}
\end{equation}
where $\Omega \subset \mathbb{R}^3$ denotes the pelvic continuum domain, $\mathbf{u}(\mathbf{x})$ is the displacement field, $\mathbf{v}$ is an arbitrary virtual displacement test function in the Sobolev space $\mathcal{V} = [H^1(\Omega)]^3$, $\boldsymbol{\varepsilon}(\mathbf{u}) = \frac{1}{2}(\nabla\mathbf{u} + (\nabla\mathbf{u})^T)$ is the infinitesimal strain tensor, and $\boldsymbol{\sigma} = \mathbf{C} : \boldsymbol{\varepsilon}$ is the Cauchy stress tensor.

Fixed anatomical boundary conditions are enforced on the superior S1 sacral facet ($\Gamma_{S1}$) via a high-stiffness Dirichlet penalty parameter:
\begin{equation}
\gamma = 10^6 \; \text{N/mm},
\end{equation}
which effectively enforces $\mathbf{u}|_{\Gamma_{S1}} \approx \mathbf{0}$ while preserving variational symmetry. External static load cases apply surface tractions $\mathbf{t}_{\text{ext}}$ integrated over specified anatomical surface facets $\Gamma_{\text{load}}$ such that $\int_{\Gamma_{\text{load}}} \mathbf{t}_{\text{ext}} \, ds = \mathbf{F}_{\text{ext}}$:
\begin{enumerate}
\item \textbf{Bilateral Stance (SP2leg):} Upward ground reaction forces applied to bilateral acetabular notches:
\begin{equation}
\mathbf{F}_{\text{AC},L} = [0, \; 0, \; +400]^T \; \text{N}, \quad \mathbf{F}_{\text{AC},R} = [0, \; 0, \; +400]^T \; \text{N} \quad (\Sigma F_{\text{CC}} = +800 \; \text{N}).
\end{equation}
\item \textbf{Unilateral Stance (SP1leg):} Upward ground reaction force applied to the right acetabular notch:
\begin{equation}
\mathbf{F}_{\text{AC},R} = [0, \; 0, \; +800]^T \; \text{N}, \quad \mathbf{F}_{\text{AC},L} = \mathbf{0}.
\end{equation}
\item \textbf{Internal Ring Compression/Shear Proxy (LAB1):} Opposing mediolateral force pair applied to the inner pectineal pelvic ring surfaces:
\begin{equation}
\mathbf{F}_{\text{ring},L} = [+400, \; 0, \; 0]^T \; \text{N}, \quad \mathbf{F}_{\text{ring},R} = [-400, \; 0, \; 0]^T \; \text{N}.
\end{equation}
\item \textbf{Ischial Distraction Proxy (LAB2):} Opposing mediolateral outward force pair applied to bilateral ischial tuberosities:
\begin{equation}
\mathbf{F}_{\text{isch},L} = [+400, \; 0, \; 0]^T \; \text{N}, \quad \mathbf{F}_{\text{isch},R} = [-400, \; 0, \; 0]^T \; \text{N}.
\end{equation}
\item \textbf{Outlet Sagittal Distraction Proxy (LAB3):} Opposing anteroposterior force pair applied to the anterior pubic insertion and posterior sacrococcygeal junction (SCJ):
\begin{equation}
\mathbf{F}_{\text{pubis}} = [0, \; -400, \; 0]^T \; \text{N}, \quad \mathbf{F}_{\text{SCJ}} = [0, \; +400, \; 0]^T \; \text{N}.
\end{equation}
\end{enumerate}
Constitutive materials comprise CT-derived heterogeneous Young's modulus $E(\mathbf{x})$ for cortical/trabecular bone ($\nu = 0.3$), isotropic articular cartilage ($E = 15$~MPa, $\nu = 0.45$), pubic symphyseal disc ($E = 5$~MPa, $\nu = 0.45$), and 1D tension-only ligament bundles with anatomically assigned stiffnesses and initial pretension.

\subsection{Construction of the anatomical pelvic coordinate frame}
To ensure objective, standardized kinematic reporting invariant to global spatial placement, an anatomical pelvic coordinate frame $\mathbf{P} \in \mathbb{R}^{3 \times 3}$ is constructed from the reference pelvic geometry $\mathcal{P}^0 = \{\mathbf{p}_i^0\}_{i=1}^N$:
\begin{enumerate}
\item \textbf{Midsagittal Symmetry Plane:} The global mediolateral unit normal $\hat{\mathbf{x}}_{\text{ML}}$ and pelvic origin $\mathbf{o} \in \mathbb{R}^3$ are determined by identifying the optimal mirror symmetry plane minimizing bilateral reflection distance across the pelvic point cloud.
\item \textbf{In-Plane Projection and Craniocaudal Axis:} Pelvic coordinates are projected onto the midsagittal symmetry plane:
\begin{equation}
\mathbf{p}_i^{\text{proj}} = (\mathbf{p}_i^0 - \mathbf{o}) - \big[(\mathbf{p}_i^0 - \mathbf{o}) \cdot \hat{\mathbf{x}}_{\text{ML}}\big] \hat{\mathbf{x}}_{\text{ML}}.
\end{equation}
Singular value decomposition (SVD) on the centered in-plane coordinates $\mathbf{X} = \mathbf{p}_i^{\text{proj}} - \bar{\mathbf{p}}^{\text{proj}} = \mathbf{U} \boldsymbol{\Sigma} \mathbf{V}^T$ extracts the dominant longitudinal principal direction $\mathbf{v}_1$. Orthogonalizing against $\hat{\mathbf{x}}_{\text{ML}}$ establishes the craniocaudal unit vector $\hat{\mathbf{z}}_{\text{CC}}$:
\begin{equation}
\hat{\mathbf{z}}_{\text{CC}} = \frac{\mathbf{v}_1 - (\mathbf{v}_1 \cdot \hat{\mathbf{x}}_{\text{ML}}) \hat{\mathbf{x}}_{\text{ML}}}{\|\mathbf{v}_1 - (\mathbf{v}_1 \cdot \hat{\mathbf{x}}_{\text{ML}}) \hat{\mathbf{x}}_{\text{ML}}\|}.
\end{equation}
\item \textbf{Anteroposterior Axis and Right-Handed System:} The anteroposterior unit vector $\hat{\mathbf{y}}_{\text{AP}}$ is established via the vector cross product:
\begin{equation}
\hat{\mathbf{y}}_{\text{AP}} = \hat{\mathbf{z}}_{\text{CC}} \times \hat{\mathbf{x}}_{\text{ML}}, \quad \text{with } \hat{\mathbf{z}}_{\text{CC}} \leftarrow \hat{\mathbf{x}}_{\text{ML}} \times \hat{\mathbf{y}}_{\text{AP}}.
\end{equation}
\item \textbf{Directional Consistency:} The sign of $\hat{\mathbf{x}}_{\text{ML}}$ is aligned such that it points from anatomical left to right: if $(\bar{\mathbf{p}}_{\text{left}} - \mathbf{o})\cdot \hat{\mathbf{x}}_{\text{ML}} > (\bar{\mathbf{p}}_{\text{right}} - \mathbf{o})\cdot \hat{\mathbf{x}}_{\text{ML}}$, then $\hat{\mathbf{x}}_{\text{ML}} \leftarrow -\hat{\mathbf{x}}_{\text{ML}}$ and $\hat{\mathbf{y}}_{\text{AP}} \leftarrow -\hat{\mathbf{y}}_{\text{AP}}$.
\end{enumerate}
The resulting orthonormal transformation matrix is:
\begin{equation}
\mathbf{P} = \begin{bmatrix} \hat{\mathbf{x}}_{\text{ML}} & \hat{\mathbf{y}}_{\text{AP}} & \hat{\mathbf{z}}_{\text{CC}} \end{bmatrix}, \quad \mathbf{P}^T \mathbf{P} = \mathbf{I}_{3 \times 3}.
\end{equation}

\subsection{Global affine sacrum drift elimination}
To isolate intra-pelvic joint mobility from any slight residual rigid-body motion or global affine deformation of the sacral base, an affine transformation $(\mathbf{A}_S, \mathbf{b}_S)$ is fitted between reference sacral coordinates $\mathbf{P}_S^0 \in \mathbb{R}^{N_S \times 3}$ and deformed sacral coordinates $\mathbf{P}_S^1 \in \mathbb{R}^{N_S \times 3}$:
\begin{equation}
\min_{\mathbf{A}_S \in \mathbb{R}^{3 \times 3}, \; \mathbf{b}_S \in \mathbb{R}^3} \sum_{i=1}^{N_S} \big\| \mathbf{p}_{S,i}^1 - (\mathbf{A}_S \mathbf{p}_{S,i}^0 + \mathbf{b}_S) \big\|^2.
\end{equation}
The optimal affine map is obtained via regularized least squares. The inverse affine transformation is then applied to all deformed bones (sacrum, left ilium, right ilium):
\begin{equation}
\mathbf{x}' = \mathbf{A}_S^{-1} (\mathbf{x}^1 - \mathbf{b}_S),
\end{equation}
anchoring the sacrum to its un-deformed reference baseline.

\subsection{Articular region of interest (ROI) identification}
For each joint side $k \in \{\text{left}, \text{right}\}$, opposing articulating bone surfaces are isolated using a proximity threshold $g = 5.0$~mm:
\begin{align}
\mathcal{R}_S^{(k)} &= \left\{ \mathbf{p} \in \mathcal{S}_{\text{sac}}^{(k)} \;\Big|\; \min_{\mathbf{q} \in \mathcal{S}_{\text{ili}}^{(k)}} \|\mathbf{p} - \mathbf{q}\| \le g \right\}, \\
\mathcal{R}_I^{(k)} &= \left\{ \mathbf{q} \in \mathcal{S}_{\text{ili}}^{(k)} \;\Big|\; \min_{\mathbf{p} \in \mathcal{S}_{\text{sac}}^{(k)}} \|\mathbf{q} - \mathbf{p}\| \le g \right\}.
\end{align}
If an extracted ROI contains fewer than $N_{\min} = 80$ nodes, the $N_{\min}$ nearest surface nodes are selected to ensure robust registration support.

\subsection{Robust trimmed Kabsch rigid registration}
For both articulating bones (sacrum $S$ and ilium $I$), rigid transformations $(\mathbf{R}_B, \mathbf{t}_B)$ mapping reference ROI points $\mathbf{P}_B^0$ to transformed deformed ROI points $\mathbf{Q}_B'$ are determined using the Kabsch algorithm with iterative outlier trimming ($\text{trim\_frac} = 0.10$, 2 iterations).

Given centered coordinates $\tilde{\mathbf{P}} = \mathbf{P}_B^0 - \bar{\mathbf{p}}_B^0$ and $\tilde{\mathbf{Q}} = \mathbf{Q}_B' - \bar{\mathbf{q}}_B'$, the cross-covariance matrix is:
\begin{equation}
\mathbf{C} = \tilde{\mathbf{P}}^T \tilde{\mathbf{Q}} \in \mathbb{R}^{3 \times 3}.
\end{equation}
Computing the singular value decomposition $\mathbf{C} = \mathbf{U} \boldsymbol{\Sigma} \mathbf{V}^T$, the optimal rotation matrix is:
\begin{equation}
\mathbf{R}_B = \mathbf{V} \begin{bmatrix} 1 & 0 & 0 \\ 0 & 1 & 0 \\ 0 & 0 & d \end{bmatrix} \mathbf{U}^T, \quad \text{where } d = \det(\mathbf{V}\mathbf{U}^T) \in \{+1, -1\},
\end{equation}
which strictly enforces $\det(\mathbf{R}_B) = +1$ and prevents improper reflection. The translation vector is:
\begin{equation}
\mathbf{t}_B = \bar{\mathbf{q}}_B' - \mathbf{R}_B \bar{\mathbf{p}}_B^0.
\end{equation}
Residual spatial deviations $r_i = \|\mathbf{R}_B \mathbf{p}_{B,i}^0 + \mathbf{t}_B - \mathbf{q}_{B,i}'\|$ are evaluated across all points; the top 10\% largest residuals are trimmed, and $(\mathbf{R}_B, \mathbf{t}_B)$ is recomputed on the remaining inliers. This prevents localized superficial element distortion from biasing macroscopic rigid-body joint estimates.

\subsection{Relative kinematics and Cardan angle decomposition}
The relative rigid transformation of the ilium with respect to the sacrum is given by:
\begin{equation}
\mathbf{R}_{\text{rel}} = \mathbf{R}_S^T \mathbf{R}_I, \quad \mathbf{t}_{\text{rel}} = \mathbf{R}_S^T (\mathbf{t}_I - \mathbf{t}_S).
\end{equation}
Projecting the relative rotation into the anatomical pelvic coordinate system yields:
\begin{equation}
\mathbf{R}_{\text{loc}} = \mathbf{P}^T \mathbf{R}_{\text{rel}} \mathbf{P} = \begin{bmatrix} R_{00} & R_{01} & R_{02} \\ R_{10} & R_{11} & R_{12} \\ R_{20} & R_{21} & R_{22} \end{bmatrix}.
\end{equation}
Factoring $\mathbf{R}_{\text{loc}}$ into an intrinsic Cardan $XYZ$ (Tait-Bryan) sequence:
\begin{equation}
\mathbf{R}_{\text{loc}} = \mathbf{R}_x(\alpha_x) \, \mathbf{R}_y(\beta_y) \, \mathbf{R}_z(\gamma_z),
\end{equation}
extracts the anatomical joint rotations:
\begin{align}
\alpha_x &= \theta_{\text{nut}} = \text{atan2}(R_{21}, \; R_{22}) \quad \text{[Nutation ($>0$) / Counternutation ($<0$) around $\hat{\mathbf{x}}_{\text{ML}}$]}, \label{eq:app_nut} \\
\beta_y &= \theta_{\text{abd}} = \arcsin(-\text{clip}(R_{20}, -1.0, 1.0)) \quad \text{[Abduction / Out-flare ($>0$) around $\hat{\mathbf{y}}_{\text{AP}}$]}, \label{eq:app_abd} \\
\gamma_z &= \theta_{\text{rot}} = \text{atan2}(R_{10}, \; R_{00}) \quad \text{[Axial spin / Torsion around $\hat{\mathbf{z}}_{\text{CC}}$]}. \label{eq:app_rot}
\end{align}
Singularities at $\beta_y = \pm 90^\circ$ ($|R_{20}| \approx 1$) are handled via $\alpha_x = \text{atan2}(-R_{01}, R_{11})$ and $\gamma_z = 0$.

\subsection{Articular contact centroid translation vector}
The anatomical reference contact location $\mathbf{p}_c$ is defined as the centroid of the reference sacral articular facet ROI:
\begin{equation}
\mathbf{p}_c = \frac{1}{|\mathcal{R}_S|} \sum_{i \in \mathcal{R}_S} \mathbf{p}_{S,i}^0.
\end{equation}
The relative displacement vector of this joint contact point in global space is:
\begin{equation}
\mathbf{d}_{\text{world}} = (\mathbf{R}_I - \mathbf{R}_S) \mathbf{p}_c + (\mathbf{t}_I - \mathbf{t}_S).
\end{equation}
Projecting $\mathbf{d}_{\text{world}}$ into the anatomical pelvic axes yields the component translations:
\begin{equation}
\mathbf{d}_{\text{pelvis}} = \mathbf{P}^T \mathbf{d}_{\text{world}} = \begin{bmatrix} d_{\text{ML}} \\ d_{\text{AP}} \\ d_{\text{CC}} \end{bmatrix}.
\end{equation}

\subsection{Bilateral kinematic aggregation and cohort summary metrics}
Relative rotations and translations are evaluated independently for the left ($L$) and right ($R$) joints. The 3D resultant norms are:
\begin{equation}
\|\boldsymbol{\theta}_k\| = \sqrt{\alpha_{x,k}^2 + \beta_{y,k}^2 + \gamma_{z,k}^2}, \quad \|\mathbf{d}_k\| = \sqrt{d_{\text{ML},k}^2 + d_{\text{AP},k}^2 + d_{\text{CC},k}^2}, \quad k \in \{L, R\}.
\end{equation}
To prevent lateral sign cancellation between symmetric left and right joints while capturing total mechanical compliance, cohort-level summary metrics are defined as bilateral means of absolute component values:
\begin{align}
\theta_{\text{mag}} &= \frac{1}{2}\big(\|\boldsymbol{\theta}_L\| + \|\boldsymbol{\theta}_R\|\big), \quad & d_{\text{mag}} &= \frac{1}{2}\big(\|\mathbf{d}_L\| + \|\mathbf{d}_R\|\big), \\
\theta_{\text{nut}} &= \frac{1}{2}\big(|\alpha_{x,L}| + |\alpha_{x,R}|\big), \quad & |d_{\text{ML}}| &= \frac{1}{2}\big(|d_{\text{ML},L}| + |d_{\text{ML},R}|\big), \\
|d_{\text{AP}}| &= \frac{1}{2}\big(|d_{\text{AP},L}| + |d_{\text{AP},R}|\big), \quad & |d_{\text{CC}}| &= \frac{1}{2}\big(|d_{\text{CC},L}| + |d_{\text{CC},R}|\big).
\end{align}
Bilateral translation asymmetry is defined as the absolute magnitude difference:
\begin{equation}
\Delta_{\text{asym}} = \Big| \|\mathbf{d}_L\| - \|\mathbf{d}_R\| \Big|.
\end{equation}

\subsection{Morphometric triad definitions}
The three structural covariates comprising the morphometric triad are evaluated directly from landmark coordinates warped to each subject's geometry:
\begin{enumerate}
\item \textbf{Anteroposterior Inlet Diameter ($d_{\text{AP}}$):} Conjugate vera, defined as the Euclidean distance between the anterior sacral promontory ($\mathbf{p}_{\text{prom}}$) and the superior/posterior pubic symphyseal crest ($\mathbf{p}_{\text{pubis}}$):
\begin{equation}
d_{\text{AP}} = \|\mathbf{p}_{\text{prom}} - \mathbf{p}_{\text{pubis}}\|.
\end{equation}
\item \textbf{Biischiadic Width ($w_{\text{biisch}}$):} Interspinous outlet diameter, defined as the transverse distance between the left and right ischial spines ($\mathbf{p}_{\text{spine},L}$ and $\mathbf{p}_{\text{spine},R}$):
\begin{equation}
w_{\text{biisch}} = \|\mathbf{p}_{\text{spine},L} - \mathbf{p}_{\text{spine},R}\|.
\end{equation}
\item \textbf{Subpubic Arch Angle ($\alpha_{\text{subpubic}}$):} Planar angle measured at the inferior apex of the pubic symphysis ($\mathbf{p}_{\text{apex}}$) between vectors directed along the descending left and right inferior pubic/ischial rami ($\mathbf{p}_{\text{rami},L}$ and $\mathbf{p}_{\text{rami},R}$):
\begin{equation}
\cos\alpha_{\text{subpubic}} = \frac{\mathbf{v}_L \cdot \mathbf{v}_R}{\|\mathbf{v}_L\| \|\mathbf{v}_R\|}, \quad \text{where } \mathbf{v}_L = \mathbf{p}_{\text{rami},L} - \mathbf{p}_{\text{apex}}, \;\; \mathbf{v}_R = \mathbf{p}_{\text{rami},R} - \mathbf{p}_{\text{apex}}.
\end{equation}
\end{enumerate}
Total pelvic bone volume ($\text{Vol}$) is evaluated via direct spatial integration over tetrahedral bone elements: $\text{Vol} = \sum_{e \in \Omega_{\text{bone}}} V_e$.

